%% ELEC2013 2007

%% ----------------------------------------------------------------
%% exam.tex
%% ----------------------------------------------------------------
\documentclass{sotonExamBoxes}    % use option [answers] to include them
%\documentclass[answers]{uosexamb}
%% ----------------------------------------------------------------
\usepackage{graphicx}
%\usepackage{pstricks}
\usepackage{alltt}
\usepackage{bm}
\usepackage{amssymb}
\usepackage{amsfonts}
\usepackage{amsmath}
\thinmuskip=1mu
%\renewcommand{\ttdefault}{pcr}
%% macros.tex -- shared macro definitions for the problem sheets.
%%
%% Every sheet inputs this instead of carrying its own copy. Input it after
%% the \usepackage lines: the definitions below use bm, amsmath and amssymb.
%%
%% Six macros were defined inconsistently across the sheets when this file
%% was made; the choices are marked CONFLICT below.




\usepackage{stmaryrd}

\renewcommand{\ttdefault}{pcr}

\newcommand{\tr}{\textsf{T}}

\newcommand{\e}[1]{{\rm e}^{#1}}

\newcommand{\E}[1]{{\rm e}^{{\displaystyle #1}}}

\newcommand{\logg}[1]{\log\left(#1\right)}

\newcommand{\Prob}[1]{\mathbb{P}\left(#1\right)}

\newcommand{\Step}{\mathop{\mathrm{O}\hspace{-15pt}\rule[8.5pt]{11pt}{1.2pt}\hspace{2pt}}}

\newcommand\diag{\mathop{\mathrm{diag}}}

\newcommand{\grad}{\bm{\nabla}}

\newcommand{\dd}{\mathrm{d}}

\newcommand{\class}{\mathcal{C}}

\newcommand{\data}{\mathcal{D}}

%% CONFLICT: gaussianProcess.tex had {U}{eur}{b}{n} (Euler bold);
%% every other sheet had this. Matrices in that one sheet now look
%% different -- see the report.

\DeclareMathAlphabet{\mat}{OT1}{cmss}{bx}{n}

%% The identity needs to be distinguishable. In cmss a bold I is a plain
%% vertical bar with no crosspieces, so \mat{I} reads as a rule rather than a
%% letter. Euler's bold I keeps its serifs, so the identity alone is set in
%% Euler while everything else stays cmss.
\DeclareMathAlphabet{\matEuler}{U}{eur}{b}{n}
\newcommand{\Id}{\matEuler{I}}

%% CONFLICT and BUG FIX: 16 sheets defined this using \qpartial,
%% which is defined nowhere, so \pd was broken in all of them.
%% Only problem2.tex had the working \partial form, kept here.

\newcommand{\pd}[2]{\ensuremath{\frac{\partial #1}{\partial #2}}\xspace}

\newcommand{\M}[1]{\ensuremath{\boldsymbol{#1}}\xspace}

\newcommand{\bst}{\rule{0mm}{5mm}}

%% CONFLICT: some sheets omitted the \! . Kept the tighter spacing,
%% which is what every sheet that actually uses \av had.

\newcommand{\av}[2][\,]{\mathbb{E}_{#1}\!\left[ {\strut #2} \right]}

\newcommand{\len}[1]{\| #1 \|}

\newcommand{\inner}[2]{\left\langle #1, #2\right\rangle}

\newcommand{\squeeze}{\setlength{\itemsep}{-2pt}}

%% CONFLICT: two spacings, identical output in maths mode.

\newcommand{\bra}[1]{\left( #1 \right)}

\newcommand{\pred}[1]{\left\llbracket { \small #1} \right\rrbracket}

\newcommand{\mln}{\hat{f}_n(\bm{x})}


%% \mathop makes "tr" an operator rather than three variables juxtaposed:
%% it is set upright and TeX adds its own spacing on either side. \nolimits
%% keeps a sub/superscript beside it rather than under it in a displayed
%% equation.
%%
%% The space is \; (thickmuskip) and not \, (thinmuskip) because every sheet
%% sets \thinmuskip=1mu, which would leave \, all but invisible. Use \: for a
%% slightly tighter gap.
\newcommand{\Tr}{\mathop{\mathrm{tr}}\nolimits\;}

%% NOT SHARED: gaussianProcess.tex redefines \det to take an argument
%% (\det{\mat{V}}), while calculus.tex and vectorSpaces.tex use the standard
%% no-argument operator (\det(\mat{X})). The two cannot coexist, so that
%% redefinition stays local to gaussianProcess.tex.

%% NOT SHARED: \ml and \mm mean different things in different sheets.
%% amgm and convexity use a no-argument \ml and \mm = \hat{m}_c(x);
%% ensembleLearning uses a one-argument \ml and \mm = \hat{m}(x).
%% Both forms are in use, so each sheet keeps its own.

\graphicspath{{figures/} {../2018-19/figures}}
\graphicspath{{figures/}}
\usepackage{parskip}
\AddToHook{env/samepage/begin}{%
\AddToHook{para/before}{\nopagebreak}%
\AddToHook{para/after}{\nopagebreak}%
}
% \usepackage{picture}
\begin{document}
\unitTitle{Calculus Problem Sheet}
\authors{Adam Pr\"ugel-Bennett}
\unitcode{AICE3001/6002}
\semester{2}
\year{2022}
\duration{45}{0}


\maketitle


% ***************************** question 1 **************************


\begin{question}{20}
  \begin{qparts}
    
      \qpart[3]{Show that if the columns, $\bm{m}_j$, of an $n\times n$ matrix,
        $\mat{M}$, are not independent then $\mat{M}$ does not have a
        unique inverse.  For the columns to not be independent implies
        there exists some set of coefficients $c_j$, not all zero, such that
        \begin{align*}
          \sum_{j=1}^n c_j\, \bm{m}_j = \bm{0}.
        \end{align*}
        We note that the $i^{th}$ element of $\bm{m}_j = M_{ij}$ so
        \begin{align*}
          \forall_{i=1}^n: \quad \sum_{j=1}^n c_j\, M_{ij} = 0.
        \end{align*}
        [Hint: consider acting on the vector $\bm{c}$ whose $i^{th}$
        element is $c_i$.  Hence show that if $\bm{x}$ is a solution
        to the equation $\bm{y} = \mat{M}\,\bm{x}$ then so is $\bm{x}
        + b\, \bm{c}$ for any $c\in \mathbb{R}$.]
      \begin{answer}
        We consider the vector $\bm{c}$ with elements $c_i$, then the
        $i^{th}$ component of $\mat{M}\,\bm{c}$ is equal to
        \begin{align*}
          [\mat{M}\,\bm{c}]_i = \sum_{j=1}^n M_{ij}\, c_j = 0
        \end{align*}
        or $\mat{M}\, \bm{c} = \bm{0}$.  Consequently for any vector
        $\bm{x}$ and any scalar $b$ we have that
        \begin{align*}
          \mat{M}\, \bra{\strut \bm{x} + b\, \bm{c}} =  \mat{M}\, \bm{x}
        \end{align*}
        if we denote $\bm{y} = \mat{M}\, \bm{x}$, then there is not a
        unique vector $\bm{x}$ such that $\bm{y} = \mat{M}\, \bm{x}$,
        but rather a family of vectors $\bm{x} + b\, \bm{c}$ for any $b$.
      \end{answer}}{\lines{5}}
    
        
      \qpart[3]{Show that for three $n\times n$ matrices $\mat{A}$,
        $\mat{B}$ and $\mat{C}$ the inverse of
        $\mat{A}\,\mat{B}\,\mat{C}$ is
        $\mat{C}^{-1}\mat{B}^{-1}\mat{A}^{-1}$.  You can do this by
        multiplying the product of matrices by its inverse and using
        associativity of matrix multiplication
      \begin{answer}
        \begin{align*}
          \bra{\mat{C}^{-1}\mat{B}^{-1}\mat{A}^{-1}} \bra{\strut
          \mat{A}\,\mat{B}\,\mat{C}}
          &=
            \mat{C}^{-1}\mat{B}^{-1}\bra{\mat{A}^{-1}\mat{A}}\,\mat{B}\,\mat{C}
          \\
          &= \mat{C}^{-1}\bra{\mat{B}^{-1}\mat{B}}\,\mat{C} =
            \mat{C}^{-1}\mat{C} = \Id
        \end{align*}
      \end{answer}}{\lines{3}}
    

    
      \qpart[3]{For non-square matrices, $\mat{A}$ and $\mat{B}$ we can
        have $\mat{A}\,\mat{B} = \Id$ while $\mat{B}\,\mat{A} \neq
        \Id$.  Show this is the case for the matrices
        \begin{align*}
          \mat{A} &=
                    \begin{pmatrix}
                      1 & 0
                    \end{pmatrix}
          & \mat{B} &=
                      \begin{pmatrix}
                        1 \\ 0
                      \end{pmatrix}
        \end{align*}
      \begin{answer}
        \begin{align*}
          \mat{A}\,\mat{B} &=
                             \begin{pmatrix}
                               1 & 0
                             \end{pmatrix}
                             \begin{pmatrix}
                               1 \\ 0
                             \end{pmatrix}
                             = (1) = \Id_1
          &
            \mat{B}\,\mat{A} &=
                               \begin{pmatrix}
                                 1 \\ 0
                               \end{pmatrix}
                               \begin{pmatrix}
                                 1 & 0
                               \end{pmatrix} 
                               =
                               \begin{pmatrix}
                                 1 & 0 \\ 0 & 0
                               \end{pmatrix}
        \end{align*}
      \end{answer}}{\lines{5}}

    
      \qpart[2]{For finite-dimensional square matrices if
        $\mat{A}\,\mat{B}= \Id$ then $\mat{B}\,\mat{A}= \Id$. Thus
        $\mat{B} = \mat{A}^{-1}$.  This is not entirely trivial to
        prove.  We break the proof down.

        Show that if $\mat{A}\mat{B}= \Id$ then $\mat{B}$ is
        injective (one-to-one).  This means that
        \begin{align}
          \mat{B}\, \bm{x} = \mat{B}\, \bm{y} \implies \bm{x} =
          \bm{y}.
          \label{eq:injective}
        \end{align}
        In other words if $\bm{x}\neq\bm{y}$ then $\mat{B}\, \bm{x}$
        must be different for $\mat{B}\, \bm{y}$.  Assuming
        $\mat{A}\mat{B}= \Id$ prove equation~(\ref{eq:injective}).
      \begin{answer}
        Assume $\mat{B}\, \bm{x} = \mat{B}\, \bm{y}$ and multiply by $\mat{A}$
        \begin{align*}
          \mat{A}\,\mat{B}\, \bm{x} &= \mat{A}\,\mat{B}\, \bm{y} \\
          \Id \bm{x} &= \Id \bm{y} \\
          \bm{x} &= \bm{y}
        \end{align*}
      \end{answer}}{\lines{8}}
    

    
      \qpart[5]{Next we want to show that in a finite dimensional space
        a square matrix that is injective is also surjective (onto).
        That is, if $\mat{B}$ is injective then for any vector $\bm{y}$
        there exists a vector $\bm{x}$ such that
        $\bm{y}= \mat{B}\bm{x}$.  The reason a finite dimensional square
        matrix that is injective is also surjective is not so obvious.
        To prove this we consider a set of $n$ \textbf{basis} vectors
        that span the entire space $\mathbb{R}^n$.  We call these
        $\bm{v}_1$, $\bm{v}_2$, \ldots, $\bm{v}_n$.  Because these are
        independent the only set of numbers $c_i$ such that
        $\sum_i c_i \bm{v}_i=\bm{0}$ is when $c_1=c_2=\ldots=c_n=0$.
        Show that the set of vectors $\bm{v}_i' = \mat{B}\, \bm{v}_i$
        form a set of basis vectors (i.e. they are linearly
        independent).  This is easily proved by contradiction.
      \begin{answer}
        Since $\mat{B}$ is injective and $\bm{v}_i\neq \bm{0}$ (since
        basis vectors are non-zero), clearly $\bm{v}_i' = \mat{B}\,
        \bm{v}_i$ are non-zero.  The non-trivial part is to show these
        vectors are linearly independent.  Suppose this is not the case and
        there exists a set of coefficients $c_i'$, not all zero, such that
        \begin{align*}
          \sum_{i=1}^n c_i'\, \bm{v}_i' &= \bm{0} \\
          \sum_{i=1}^n c_i'\, \mat{B} \, \bm{v}_i &= \bm{0} \\
          \mat{B} \, \sum_{i=1}^n c_i'\, \bm{v}_i &= \bm{0}.      
        \end{align*}
        But, as $\mat{B}$ is injective this implies
        \begin{align*}
          \sum_{i=1}^n c_i'\, \bm{v}_i &= \bm{0}.
        \end{align*}
        Since the vectors $\bm{v}_i$ are linearly independent this
        implies every $c_i'=0$, contradicting our assumption.
      \end{answer}}{\lines{8}}
    

    
      \qpart[2]{Use the fact that the vectors $\bm{v}_i'$ are linearly
        independent to show that, for any vector $\bm{y}$, we can find a
        vector $\bm{x}$ such that $\bm{y} = \mat{B}\,\bm{x}$
      \begin{answer}
        Any vector $\bm{y}$ can be represented by a weighted sum of
        $n$ independent vectors.  In particular we can write
        \begin{align*}
          \bm{y} &= \sum_{i=1}^n c_i\, \bm{v}'_i = \sum_{i=1}^n c_i\,
                   \mat{B} \, \bm{v}_i  \\
                 &= \mat{B} \bra{\sum_{i=1}^n c_i\, \bm{v}_i} = \mat{B}\,\bm{x}
        \end{align*}
        where $\bm{x} = \sum_{i=1}^n c_i\, \bm{v}_i$.
      \end{answer}}{\lines{6}}
    

    
      \qpart[2]{Using the fact that $\mat{A}\,\mat{B}= \Id$ and
        $\mat{B}$ is surjective show $\mat{B}\,\mat{A} =
        \Id$
      \begin{answer}
        We note that for all $\bm{x}$
        \begin{align*}
          \mat{A}\,\mat{B}\,\bm{x} &= \bm{x} \\
          \mat{B}\,\mat{A}\,\mat{B}\,\bm{x} &= \mat{B}\,\bm{x} 
        \end{align*}
        But as $\mat{B}$ is surjective for any vector $\bm{y}$ there
        exists a vector $\bm{x}$ such that $\bm{y} = \mat{B}\,\bm{x}$.
        Thus,
        \begin{align*}
          \mat{B}\,\mat{A}\,\bm{y} = \bm{y},
        \end{align*}
        but as this is true for any vector $\bm{y}$ it follows that $\mat{B}\,\mat{A}=\Id$.
      \end{answer}}{\lines{5}}
    
  \end{qparts}
\end{question}
\freshpage

% ***************************** question 2 **************************

\begin{question}{20}
  Throughout this question you may use two standard properties of the
  trace, namely that a matrix and its transpose have the same trace, and
  that the trace is unchanged by cyclically permuting a product
  \begin{align*}
    \Tr \mat{A} &= \Tr \mat{A}^\tr
    & \Tr \mat{A}\,\mat{B} &= \Tr \mat{B}\,\mat{A} .
  \end{align*}

  \begin{qparts}

    \qpart[5]{For a vector $\bm{a}\in\mathbb{R}^m$, a vector
      $\bm{b}\in\mathbb{R}^n$ and an $m\times n$ matrix $\mat{U}$, show
      that
      \begin{align*}
        \bm{a}^\tr\, \mat{U}\, \bm{b}
        = \Tr\bra{\strut \mat{U}^\tr\, \bm{a}\, \bm{b}^\tr}.
      \end{align*}
      We can show this by writing out this in terms of components (as
      we did in the lectures), however, we can do this much faster by
      treating $\bm{a}^\tr\, \mat{U}\, \bm{b}$ as a $1\times1$ matrix
      which is equal to it trace: $\Tr \bm{a}^\tr\, \mat{U}\, \bm{b}$.
      \begin{answer}
        Thus
        \begin{align*}
          \bm{a}^\tr\, \mat{U}\, \bm{b}
          &= \Tr \bm{a}^\tr\,\mat{U}\, \bm{b} \\
          &= \Tr \bra{\bm{a}^\tr\,\mat{U}\, \bm{b}}^\tr
            =  \Tr \bm{b}^\tr\,\mat{U}^\tr\, \bm{a} \\
            &= \Tr \mat{U}^\tr\, \bm{a} \, \bm{b}^\tr
        \end{align*}
      \end{answer}}{\lines{6}}

    \qpart[4]{Verify the identity explicitly for
      \begin{align*}
        \bm{a} &= \begin{pmatrix} 1 \\ 2 \end{pmatrix}
        & \bm{b} &= \begin{pmatrix} 3 \\ 1 \end{pmatrix}
        & \mat{U} &= \begin{pmatrix} 1 & 2 \\ 3 & 4 \end{pmatrix}
      \end{align*}
      by computing both $\bm{a}^\tr\,\mat{U}\,\bm{b}$ and
      $\Tr\bra{\strut\mat{U}^\tr\,\bm{a}\,\bm{b}^\tr}$.
      \begin{answer}
        For the left-hand side,
        \begin{align*}
          \mat{U}\,\bm{b} &=
          \begin{pmatrix} 1 & 2 \\ 3 & 4 \end{pmatrix}
          \begin{pmatrix} 3 \\ 1 \end{pmatrix}
          = \begin{pmatrix} 5 \\ 13 \end{pmatrix}
          & \bm{a}^\tr\,\mat{U}\,\bm{b} &=
          \begin{pmatrix} 1 & 2 \end{pmatrix}
          \begin{pmatrix} 5 \\ 13 \end{pmatrix}
          = 5 + 26 = 31 .
        \end{align*}
        For the right-hand side, the outer product is
        \begin{align*}
          \bm{a}\,\bm{b}^\tr =
          \begin{pmatrix} 1 \\ 2 \end{pmatrix}
          \begin{pmatrix} 3 & 1 \end{pmatrix}
          = \begin{pmatrix} 3 & 1 \\ 6 & 2 \end{pmatrix}
        \end{align*}
        and hence
        \begin{align*}
          \mat{U}^\tr\, \bm{a}\, \bm{b}^\tr =
          \begin{pmatrix} 1 & 3 \\ 2 & 4 \end{pmatrix}
          \begin{pmatrix} 3 & 1 \\ 6 & 2 \end{pmatrix}
          = \begin{pmatrix} 21 & 7 \\ 30 & 10 \end{pmatrix}
        \end{align*}
        whose trace is $21 + 10 = 31$, agreeing with the left-hand side.

        Note that only the diagonal elements are needed, so the
        off-diagonal entries $7$ and $30$ need never be computed.  Note
        also that $\mat{U}$ is not symmetric: had the transpose been
        omitted the two sides would not have agreed, which is the point
        of checking with a non-symmetric matrix.
      \end{answer}}{\lines{10}}


    \qpart[5]{For a square matrix $\mat{U}$ and a sufficiently small
      scalar $\epsilon$, the inverse of $\Id + \epsilon\, \mat{U}$ is
      given by the geometric (Neumann) series
      \begin{align*}
        \bra{\Id + \epsilon\, \mat{U}}^{-1}
        = \Id - \epsilon\, \mat{U} + \epsilon^2 \mat{U}^2
          - \epsilon^3 \mat{U}^3 + \cdots .
      \end{align*}
      Show this by multiplying $\Id + \epsilon\, \mat{U}$ by the
      truncated series
      \begin{align*}
        \mat{S}_n = \sum_{k=0}^{n} \bra{-\epsilon}^k \mat{U}^k
      \end{align*}
      and showing that the result differs from $\Id$ only at order
      $\epsilon^{n+1}$.  Hence write down
      $\bra{\Id + \epsilon\, \mat{U}}^{-1}$ to first order in
      $\epsilon$.
      \begin{answer}
        Multiplying out, and using
        $\epsilon\, \bra{-\epsilon}^k = -\bra{-\epsilon}^{k+1}$,
        \begin{align*}
          \bra{\Id + \epsilon\, \mat{U}}\mat{S}_n
          &= \sum_{k=0}^{n} \bra{-\epsilon}^k \mat{U}^k
             + \epsilon\, \mat{U} \sum_{k=0}^{n} \bra{-\epsilon}^k \mat{U}^k
          \\
          &= \sum_{k=0}^{n} \bra{-\epsilon}^k \mat{U}^k
             - \sum_{k=0}^{n} \bra{-\epsilon}^{k+1} \mat{U}^{k+1}
          \\
          &= \sum_{k=0}^{n} \bra{-\epsilon}^k \mat{U}^k
             - \sum_{j=1}^{n+1} \bra{-\epsilon}^{j} \mat{U}^{j} .
        \end{align*}
        The two sums are identical except for the $k=0$ term of the first
        and the $j=n+1$ term of the second, so everything between
        telescopes and
        \begin{align*}
          \bra{\Id + \epsilon\, \mat{U}}\mat{S}_n
          = \Id - \bra{-\epsilon}^{n+1} \mat{U}^{n+1} .
        \end{align*}
        The correction is of order $\epsilon^{n+1}$, so $\mat{S}_n$ is the
        inverse up to that order, and taking $n\to\infty$ gives the full
        series provided $\epsilon$ is small enough that
        $\epsilon^{n} \mat{U}^{n} \to \bm{0}$ (it suffices that
        $|\epsilon|$ is smaller than the reciprocal of the largest
        magnitude eigenvalue of $\mat{U}$).

        Keeping only the first two terms,
        \begin{align*}
          \bra{\Id + \epsilon\, \mat{U}}^{-1}
          = \Id - \epsilon\, \mat{U} + \mathrm{O}\bra{\epsilon^2} .
        \end{align*}
      \end{answer}}{\lines{10}}

    \qpart[6]{Use this to find the derivative of
      $g(\mat{X}) = \bm{a}^\tr\, \mat{X}^{-1}\, \bm{b}$ with respect to
      the matrix $\mat{X}$.  Replace $\mat{X}$ by
      $\mat{X} + \epsilon\, \mat{U}$, expand to first order in $\epsilon$,
      and use the definition of the derivative with respect to a matrix
      \begin{align*}
        \lim_{\epsilon\rightarrow 0}
        \frac{g(\mat{X} + \epsilon\, \mat{U}) - g(\mat{X})}{\epsilon}
        = \Tr\bra{\mat{U}^\tr\,
          \frac{\partial g(\mat{X})}{\partial \mat{X}}}
      \end{align*}
      which holds for every $\mat{U}$.  Part~(a) is what lets you read the
      derivative off.
      \begin{answer}
        Factorising $\mat{X}$ out on the left,
        $\mat{X} + \epsilon\, \mat{U}
        = \mat{X}\bra{\Id + \epsilon\, \mat{X}^{-1}\mat{U}}$, so by
        the previous part
        \begin{align*}
          \bra{\mat{X} + \epsilon\, \mat{U}}^{-1}
          &= \bra{\Id + \epsilon\, \mat{X}^{-1}\mat{U}}^{-1}\mat{X}^{-1}
          = \bra{\Id - \epsilon\, \mat{X}^{-1}\mat{U}
            + \mathrm{O}\bra{\epsilon^2}}\mat{X}^{-1}
          \\
          &= \mat{X}^{-1} - \epsilon\, \mat{X}^{-1}\mat{U}\,\mat{X}^{-1}
            + \mathrm{O}\bra{\epsilon^2} .
        \end{align*}
        Therefore
        \begin{align*}
          g(\mat{X} + \epsilon\, \mat{U}) - g(\mat{X})
          = -\epsilon\, \bm{a}^\tr \mat{X}^{-1}\mat{U}\,\mat{X}^{-1} \bm{b}
            + \mathrm{O}\bra{\epsilon^2}
        \end{align*}
        and dividing by $\epsilon$ and letting $\epsilon\rightarrow 0$
        \begin{align*}
          \lim_{\epsilon\rightarrow 0}
          \frac{g(\mat{X} + \epsilon\, \mat{U}) - g(\mat{X})}{\epsilon}
          = -\bm{a}^\tr \mat{X}^{-1}\, \mat{U}\, \mat{X}^{-1} \bm{b} .
        \end{align*}
        This is of the form $\bm{c}^\tr\, \mat{U}\, \bm{d}$ with
        $\bm{c}^\tr = -\bm{a}^\tr\mat{X}^{-1}$, that is
        $\bm{c} = -\bra{\mat{X}^{-1}}^\tr \bm{a}$, and
        $\bm{d} = \mat{X}^{-1}\bm{b}$.  By part~(a),
        \begin{align*}
          \bm{c}^\tr\, \mat{U}\, \bm{d}
          = \Tr\bra{\strut \mat{U}^\tr\, \bm{c}\, \bm{d}^\tr} .
        \end{align*}
        Comparing with the definition above, and since this holds for
        every $\mat{U}$, the derivative is
        \begin{align*}
          \frac{\partial g(\mat{X})}{\partial \mat{X}}
          = \bm{c}\,\bm{d}^\tr
          = -\bra{\mat{X}^{-1}}^\tr \bm{a}\, \bm{b}^\tr
             \bra{\mat{X}^{-1}}^\tr .
        \end{align*}
        As a check on the shapes, $\bm{a}\,\bm{b}^\tr$ is a matrix of the
        same size as $\mat{X}$, as the derivative must be.
      \end{answer}}{\lines{14}}

  \end{qparts}
\end{question}


% ***************************** question 3 **************************

\freshpage
\begin{question}{20}
  We want to compute the derivative of $\logg{\det(\mat{X})}$ with
  respect to the matrix $\mat{X}$.  Our strategy will be to use the
  fact that
  \begin{align*}
    \lim_{\epsilon\rightarrow 0} \frac{\logg{\det(\mat{X} + \epsilon\, \mat{U})} -
    \logg{\det\bra{\mat{X}}}}{\epsilon} = \Tr  \mat{U}^\tr \frac{\partial
    \logg{\det(\mat{X})}}{\partial \mat{X}}.
  \end{align*}
  We therefore expand $\logg{\det(\mat{X}+ \epsilon\, \mat{U})}$ to
  first order in $\epsilon$, subtract $\logg{\det\bra{\mat{X}}}$, divide by
  $\epsilon$, rearrange and read off $\frac{\partial
    \logg{\det(\mat{X})}}{\partial \mat{X}}$.  However, rearranging in
  this case is quite complicated.
  
  \begin{qparts}
    \qpart[5]{Using $\det(\mat{A}\,\mat{B}) =
      \det(\mat{A})\,\det(\mat{B})$ and $\log(a\,b)= \log(a) +
      \log(b)$ expand $\logg{\det(\mat{X}+ \epsilon \mat{U})}$ as a
      sum of $\logg{\det(\mat{X})}$ and another logarithm.
      \begin{answer}
        \begin{align*}
          \logg{\det\bra{\strut \mat{X}+\epsilon\,\mat{U}}}
          &= \logg{\det\bra{\mat{X}\,\bra{\Id + \epsilon\,\mat{X}^{-1}\mat{U}}}}\\
          &= \logg{\det(\mat{X}) \, \det(\Id + \epsilon\,\mat{X}^{-1}\mat{U})}\\
          &= \logg{\det(\mat{X})} + \logg{\det\bra{\Id +
            \epsilon\,\mat{X}^{-1}\mat{U}}}
        \end{align*}
      \end{answer}}{\lines{5}}

    \qpart[2]{The hard part is to expand $\det\bra{\Id +
        \epsilon\,\mat{M}}$. Write out $\det\bra{\Id +
        \epsilon\,\mat{M}}$ explicitly for a 2 
      dimensional matrix and expand to first order in $\epsilon$
      \begin{answer}
        \begin{align*}
          \det\bra{\Id+\epsilon\,\mat{M}}
          &=
            \det\begin{pmatrix}
              1+\epsilon\, M_{11} & \epsilon\, M_{12}\\
              \epsilon\, M_{21} & 1+\epsilon\, M_{22}
            \end{pmatrix}
            = (1+\epsilon\, M_{11})\,(1+\epsilon\, M_{22}) -
            \epsilon^2\, M_{21}\,M_{12}\\
          &= 1 + \epsilon \bra{M_{11} + M_{22}} + O(\epsilon^2)
        \end{align*}
      \end{answer}}{\lines{8}}

    \qpart[5]{For an $n\times n$ matrix $\mat{M}$ compute the
      determinant of
      \begin{align*}
        \det\bra{\Id + \epsilon\, \mat{M}} = \det
        \begin{pmatrix}
          1 + \epsilon\, M_{11} & \epsilon\, M_{12} & \ldots & \epsilon\,M_{1n}  \\
          \epsilon \, M_{21} & 1 + \epsilon\, M_{22} & \ldots &\epsilon\,M_{2n}\\
          \vdots & \vdots & \ddots & \vdots \\
          \epsilon \, M_{n,1} &  \epsilon\, M_{n2} & \ldots &1+\epsilon\,M_{nn}
        \end{pmatrix}
      \end{align*}
      up to first order in $\epsilon$.  Recall we can compute the
      determinant of a matrix $\mat{X}$ as
      \begin{align*}
        \det(\mat{X}) = \sum_{j=1}^n (-1)^{1+j}\, X_{1j} \, \det(\hat{\mat{X}}_{1j})
      \end{align*}
      where $\hat{\mat{X}}_{1j}$ is the matrix obtained
      by removing row 1 and column $j$ from the matrix $\mat{X}$.
      \begin{answer}
        Expanding the determinant we get a term equal to the
        product of the diagonal elements
        \begin{align*}
          \prod_{i=1}^n (1+\epsilon \, M_{ii}) = 1 +\epsilon\,
          \sum_{i=1}^n M_{ii} + O(\epsilon^2)
        \end{align*}
        all other terms in the determinant involve terms of at least
        order $\epsilon^2$ so do not contribute to order $\epsilon$.
        In consequence
        \begin{align*}
          \det\bra{\Id + \epsilon\, \mat{M}} = 1 + \epsilon
          \Tr  \mat{M} + O(\epsilon^2)
        \end{align*}
      \end{answer}}{\lines{8}}

    \qpart[3]{For a more elegant proof we can use the fact that
      \begin{align*}
        \det(\mat{M}) &= \prod_{i=1}^n \lambda_i& \Tr\mat{M} &=
                                                                \sum_{i=1}^n \lambda_i
      \end{align*}
      where $\lambda_i$ are the eigenvalues of the matrix $\mat{M}$.
      We will prove these results later when we look at eigen-systems.
      Note that if $\bm{v}_i$ is an eigenvector of the matrix
      $\mat{M}$, with eigenvalue $\lambda_i$ then it is an eigenvector
      of $\Id + c\, \mat{M}$ with eigenvalue $1+c\,\lambda_i$
      since
      \begin{align*}
        (\Id + c\, \mat{M})\, \bm{v}_i = \bm{v}_i + c\,\lambda_i\,
        \bm{v}_i
        = (1+c\,\lambda_i)\, \bm{v}_i.
      \end{align*}
      Using these results show that $\det\bra{\Id + \epsilon\, \mat{M}} = 1 + \epsilon\,
      \Tr  \mat{M} + O(\epsilon^2)$.
      \begin{answer}
        \begin{align*}
          \det\bra{\Id + \epsilon\, \mat{M}}
          &= \prod_{i=1}^n (1+\epsilon\, \lambda_i)
            = 1 + \epsilon\, \sum_{i=1}^n \lambda_i + O(\epsilon^2) \\
          &= 1 + \epsilon\, \Tr \mat{M} + O(\epsilon^2)
        \end{align*}
      \end{answer}}{\lines{5}}

    \qpart[5]{Using the results
      \begin{align*}
        \logg{\det(\mat{X} + \epsilon\, \mat{U})} - \logg{\det\bra{\mat{X}}}
        &= \logg{\det\bra{\Id + \epsilon\,\mat{X}^{-1}\mat{U}}} \\
        \det\bra{\Id + \epsilon\, \mat{M}} &= 1 + \epsilon
      \Tr  \mat{M} + O(\epsilon^2)
      \end{align*}
      compute $\frac{\partial \logg{\det(\mat{X})}}{\partial
        \mat{X}}$.  You may use the fact that $\log(1+x) = x +
      O(x^2)$ and $\Tr \mat{A} = \Tr \mat{A}^\tr$.
      \begin{answer}
        We note that
        \begin{align*}
          \logg{\det(\mat{X} + \epsilon\, \mat{U})} - \logg{\det\bra{\mat{X}}}
          &= \logg{1 + \epsilon\, \Tr  \mat{X}^{-1}\mat{U} +
          O(\epsilon^2)} \\
          &= \epsilon\,  \Tr  \mat{X}^{-1}\mat{U} +
          O(\epsilon^2)  \\
          &= \epsilon\, \Tr  \mat{U}^\tr (\mat{X}^{-1})^\tr
            + O(\epsilon^2)
        \end{align*}
        Thus,
        \begin{align*}
          \frac{\partial \logg{\det(\mat{X})}}{\partial \mat{X}} =  (\mat{X}^{-1})^\tr
        \end{align*}
      \end{answer}}{\lines{7}}
  \end{qparts}
\end{question}
\freshpage

% ***************************** question 4 **************************
\begin{question}{20}
    In this question I want you to use a numerical package to do some
    simple algebra.  You can use MATLAB, OCTAVE, numpy in Python or
    any other language of your choosing.  The answers are only going
    to be accurate up to the numerical precision of your machine
    (typically around $10^{-16}$).  For python the following may be
    useful
    \begin{tabularx}{\linewidth}{@{}l|X@{}}
      \hline
      \verb+import numpy as np+ & import the numpy module\\ \hline
    \verb+A = np.random.randn(n,m)+ & generate a matrix of
      normally distributed random deviates with n rows and m columns\\ \hline
    \verb+A.transpose()+ & return the transpose of $\mat{A}$
      (does not change $\mat{A}$)\\ \hline
    \verb+A@B+ & matrix multiplication $\mat{A}\mat{B}$.  Beware
      that \verb+A*B+ performs elementwise multiplication which is
      different to normal matrix multiplication\\ \hline
    \verb+A.trace()+ & returns the trace of a matrix\\ \hline
    \verb+np.linalg.det(A)+ & return the determinant of a matrix\\ \hline
    \verb+np.round(A, 10)+ & rounds the elements in a matrix
      $\mat{A}$ to 10 decimal places.  Makes expressions like
      \verb+(A@B)@C - A@(B@C)+ easier to read\\ \hline   
    \verb+np.linalg.inv(A)+ & inverse of matrix $\mat{A}$\\ \hline   
    \end{tabularx}

    Generate three $5\times 5$ random matrices $\mat{A}$, $\mat{B}$ and
    $\mat{C}$ and two random vectors of length 5, $\bm{a}$ and $\bm{b}$.

   \begin{qparts}
   \qpart[2]{Compute $\mat{D}= \mat{A}\, \mat{B}$ and $\mat{E} =
     \mat{B} \, \mat{C}$.  Which quantities are the same (a)
     $\mat{A}\,\mat{B}\, \mat{C}$, (b) $\mat{D}\, \mat{E}$, (c)
     $\mat{D}\, \mat{C}$, (d) $\mat{A}\,\mat{E}$, (e) $\mat{C}\,\mat{D}$
      \begin{answer}
        (a) = (c) = (d)
      \end{answer}}{\lines{1}}

    \qpart[2]{Compute $\mat{A}\, \mat{B}$ and $\mat{B} \, \mat{A}$.
      Are these the same?
      \begin{answer}
        Almost certainly these are different.
      \end{answer}}{\lines{1}}
    
    \qpart[2]{Compute (a) $\bm{a}^\tr \mat{A} \, \bm{b}$, (b)
      $\bm{a}^\tr \mat{A}^\tr \bm{b}$ (c) $\bm{b}^\tr \mat{A}^\tr
      \bm{a}$ and (d) $\bm{b}^\tr \mat{A}  \bm{a}$.  Which quantities
      are the same?
      \begin{answer}
        (a) = (c) and (b) = (d)
      \end{answer}}{\lines{1}}

    \qpart[2]{
      Compute (a) $\Tr \mat{A}\, \mat{B}$, (b)
      $\bra{\Tr \mat{A}}\bra{\Tr \mat{B}}$ (c)
        $\Tr \mat{B}\, \mat{A}$ and (d) $\Tr
        \mat{B}^\tr\, \mat{A}^\tr$.  Which quantities are the same?
        \begin{answer}
          (a) = (c) = (d)
      \end{answer}}{\lines{1}}

    \qpart[2]{
       Compute (a) $\det(\mat{A}\, \mat{B})$, (b)
      $\bra{\det \mat{A}}\bra{\det \mat{B}}$ (c)
        $\det (\mat{B}\, \mat{A})$ and (d) $\det(
        \mat{B}^\tr\, \mat{A}^\tr)$.  Which quantities are the same?
        \begin{answer}
          These are all the same.
      \end{answer}}{\lines{1}}

    \qpart[10]{
      Although we can define the gradient with respect to a
      matrix using
      \begin{align*}
      \lim_{\epsilon\rightarrow 0}  \frac{f(\mat{A} +
        \epsilon\,\mat{U}) - f(\mat{A})}{\epsilon} = \Tr
        \mat{U}^\tr \frac{\partial f(\mat{A})}{\partial \mat{A}}
      \end{align*}
      numerically the convergence is slow.  This is problematic as $f(\mat{A} +
      \epsilon\,\mat{U})$ and $f(\mat{A})$ are very close in value
      so when we take their difference we lose a lot of precision.
      A ratio that converges much more quickly is
      \begin{align*}
        \lim_{\epsilon\rightarrow 0}
        \frac{f(\mat{A} + \epsilon\,\mat{U}) -
        f(\mat{A}-\epsilon\,\mat{U})}{2\,\epsilon}
        = \Tr \mat{U}^\tr
        \frac{\partial f(\mat{A})}{\partial \mat{A}}.       
      \end{align*}
      If we want to obtain a more accurate numerical estimate of a
      derivative we would compute the ratio for different values of
      $\epsilon$ and then extrapolate this series to when
      $\epsilon\rightarrow0$.  However, for this exercise this is
      unnecessary.   Let $f(\mat{A}) = \logg{\det(\mat{A})}$, set
      $\mat{U} = \mat{B}$ and compute
      \begin{align*}
        \Tr \mat{U}^\tr \frac{\partial f(\mat{A})}{\partial \mat{A}}
      \end{align*}
      numerically.  Compare this to
      $\Tr \mat{U}^\tr (\mat{A}^{-1})^\tr$
      \begin{answer}
        Use the following
\begin{alltt}
eps = 0.0001\\
estimate = (np.log(np.abs(np.linalg.det(A+eps*B)))\\
\hspace*{2em}- np.log(np.abs(np.linalg.det(A-eps*B)))) / (2*eps)\\
actual = (B.transpose()\\
\hspace*{2em}@(np.linalg.inv(A)).transpose()).trace()
\end{alltt}
      \end{answer}}{\lines{5}}

  \end{qparts}
\end{question}


% ***************************** End of exam *****************************
\end{document}




